\newpage
\appendix
% \onecolumn


\begin{center}
\Large
\textbf{Appendix}
\end{center}



In the Appendix, we introduce more details  along with additional experimental results, discussions, and related works:

\begin{itemize}
\item Appendix \ref{app_sec:implem_details}: Experimental setups (cf. Section \ref{sec:exp}).
\item Appendix \ref{app_sec:more_results}: More experimental results (cf. Section \ref{sec:method} and \ref{sec:exp}).
\item Appendix \ref{app_sec:proof}: Proof of the Theorem~\ref{theorem:subspace_overlap} (cf. Section \ref{sec:method}).
\item Appendix \ref{app_sec:more_related_works}: Additional discussions and more related works (cf. Section \ref{sec:related_works}).
\end{itemize}




\captionsetup[subtable]{position=top}
\def\zsrex{Which continent is Berkner Island in?}
\def\zsrey{South America}
\def\zsrelocx{who gets the golden boot if its a tie?}
\def\zsrelocy{shared}
\def\zsrexx{On which continent is Berkner Island located?}
\def\zsreyy{South America}

\def\hallucx{This is a Wikipedia passage about heinz christian pander. Heinz Christian Pander (1794 - 1865) was a German anatomist and embryologist who was born in Riga, Latvia. He studied medicine at the University of Dorpat and later at the University of Berlin.}
\def\hallucy{In 1820, he took part in a scientific expedition to Bokhara as a naturalist.}
\def\halluclocx{Tired and restlessly, drifting in and out of sleep. Hearing crashing and banging, thinking the roof will cave in. Not alert enough to quite know what}
\def\halluclocy{it was, I yelled loudly for whoever was making those noises at such an hour to stop. They heard and listened, I'm guessing}

\def\temporalx{Self-driving cars,}
\def\temporaly{also known as autonomous vehicles, are vehicles that are capable of navigating and operating without human intervention. These innovative vehicles rely on a combination of advanced sensors, artificial intelligence, and computer algorithms to interpret their environment and make real-time decisions. With the potential to significantly impact numerous industries and sectors, self-driving cars have the ability to revolutionize transportation by enhancing safety, improving traffic flow, and increasing energy efficiency. However, challenges related to regulatory frameworks, ethical considerations, and public acceptance still need to be addressed before widespread adoption becomes a reality.}
\def\temporallocx{Apple has a new peach with the release of its 3.0GHz, 8-core Intel Xeon-based Mac Pro. The 8-core Mac Pro is powered bu two quad-core Intel Xeon \"Clov}
\def\temporallocy{ertown\" processors running at 3.0GHz. Apple also released a quad-core Mac Pro featuring two Dual-Core Intel Xeon \"Woodcrest\" processors.}
\def\temporalyy{also known as autonomous cars or driverless cars, are vehicles capable of traveling without human input. These cars utilize a range of sensors, including optical and thermographic cameras, radar, lidar, ultrasound/sonar, GPS, odometry, and inertial measurement units, to perceive their surroundings. By interpreting sensory information, control systems in the car are able to create a three-dimensional model of its environment. Using this model, the car can then identify the best navigation path and develop strategies for managing traffic controls and obstacles. As self-driving car technology continues to advance, it is expected to have a significant impact on various fields such as the automotive industry, health, welfare, urban planning, traffic, insurance, and the labor market. The regulation of autonomous vehicles is also becoming an increasingly important topic of discussion.}

\section{Implementation Details}\label{app_sec:implem_details}
\subsection{Description of Datasets}\label{app_subsec:data_description}


\begin{table}[!htbp]
\vspace{-11pt}
    \caption{Bolded text refers to the edit labels $\mathbf{y}_e$. Locality example $\mathbf{x}_{\text{loc}}$ is an unrelated query.}
    \subfloat[\textbf{ZsRE, question-answering} editing dataset example. ]{
        \centering
        \begin{tabular}{l@{\hspace{5pt}}p{2.8cm}}
            \toprule
            $\mathbf{x}_e, \mathbf{y}_e$ & \zsrex~\textbf{\zsrey} \\
            \cmidrule{2-2}
            $\mathbf{x}_{\text{loc}}$ & \zsrelocx ~ \textbf{\zsrelocy} \\
            \cmidrule{2-2}
            $\mathbf{x}'_e, \mathbf{y}_e$ & \zsrexx~\textbf{\zsreyy} \\
            \bottomrule
            \\
            \\
        \end{tabular}
    }\hfill
    \subfloat[\textbf{Hallucination} editing dataset example. In the original data \cite{grace}, there is no paraphrase $x_{e'}$ so the measurement of Gen. metric is ignored here.]{
        \centering
        \begin{tabular}{l@{\hspace{5pt}}p{7.9cm}}
            \toprule
            $\mathbf{x}_e, \mathbf{y}_e$ & \hallucx~\textbf{\hallucy} \\
            \cmidrule{2-2}
            $\mathbf{x}_{\text{loc}}$ & \halluclocx ~ \halluclocy \\
            \bottomrule
        \end{tabular}
    }
    \label{tab:dataset-example}
    \vspace{-13pt}
\end{table}


\paragraph{ZsRE}
The ZsRE question-answering task \cite{zsre} is extensively studied within the model editing literature \cite{rome, memit, mend, malmen, t-patcher}, where each record contains an editing statement $\mathbf{x}_e$, a paraphrase prompt $\mathbf{x}'_e$, and a locality prompt $\mathbf{x}_{\text{loc}}$.
We use the same train/test split as \cite{mend} (163196/19086).
Notably, only MEND requires fitting a hypernetwork on the training set; other methods discard the training set and perform edits and evaluations on the test set. 
In practice, we randomly sample 1K and 3K records from the test set to form the edit sets in Section \ref{subsec: main_results} and \ref{subsec: scaling2_3K}.

\paragraph{Hallucination}
\begin{wrapfigure}{r}{0.4\columnwidth}
    \centering
    \vspace{-13pt}
    \includegraphics[width=0.4\textwidth]{img/halluc_length_range.pdf}
    \vspace{-17pt}%
    \caption{Hallucination length statistics.}
    \label{fig:halluc_length}
    \vspace{-13pt}%
\end{wrapfigure}
We utilize the same dataset as GRACE, SelfCheckGPT \cite{selfcheckgpt}, to assess the ability of Model Editors to mitigate hallucinations in autoregressive LMs.
This setting involves editing highly inaccurate sentences (sourced from GPT-3 \cite{gpt-3}) and replacing them with corresponding sentences from actual Wikipedia entries.
This dataset aligns more closely with real-world deployment scenarios where models trigger "unexpected behaviors," and the token length of edits is significantly longer than in past datasets, making it a more challenging editing setting.
Unlike GRACE, which used GPT2-XL (1.5B) \cite{gpt-2}, our main experiments deploy larger LLMs, LLaMA and Mistral, both with 7B parameters, we measure retention of pretraining data ($\mathbf{x}_{\text{loc}}$) from the base model: RedPajama \cite{redpajama}, a public version of LLaMA’s pretraining data.
Some of the exceptionally long editing samples cannot even be accommodated on an NVIDIA A800 (80GB) due to resource limitations.
As shown in Figure \ref{fig:halluc_length}, the original dataset provided by GRACE, after tokenization with \textsc{LlamaTokenizer}, has length distributions ranging from [17,390].
The dimension of a single MLP layer in \texttt{llama-2-7b-hf} is (11008, 4096) \footnote{\url{https://huggingface.co/meta-llama/Llama-2-7b-hf}}.
Theoretically, fine-tuning an input of length 390 with default full precision and the Adam optimizer would require (390+4+4+4) * (11008 * 4096 * 4) + 4 * 7B = \textbf{100.36GB} of VRAM (for activations, gradients, first-order, and second-order optimizers), exceeding the memory capacity of the NVIDIA A800.
Consequently, we excluded excessively long samples (limiting tokenized lengths to 254) and ultimately retained 906 editing instances (compared to 1392 in GRACE).
To facilitate a fair comparison with MEND, we specifically allocated a training set for MEND, with a final train/test split of 306/600.
All methods were edited and evaluated on the test set.\looseness=-1

\paragraph{Temporal}
\cite{canonical_examples} sources the prefix $\mathbf{x}_e$ from the first paragraph of an entity's Wikipedia page and samples a paragraph $\mathbf{y}_e$ discussed by GPT-4 \cite{gpt4} about the emerging entity $\mathbf{x}_e$, which is usually noisy but may contain helpful information.
These are presented as editing prompts to Model Editors.
For out-of-distribution (OOD) generalization to complex natural contexts (not fitted), $\mathbf{y}_{\text{ood}}$ is taken from the actual Wikipedia suffix of $\mathbf{x}_e$.
This setup is utilized to evaluate the OOD generalization of Model Editors centered around a single canonical example.
Consistent with previous work \cite{grace}, the out-of-scope data $\mathbf{x}_{\text{loc}}$ is derived from the Pile \cite{pile}, the pretraining corpus of GPT-J-6B.
Examples from the dataset can be seen in Table \ref{tab:temporal-example}. 
To measure the OOD generalization of editing methods for emerging entities, we perform model editing using standardized simple examples and then evaluate this behavior on more complex instances.
Following \cite{canonical_examples}, in a natural setting, no single correct continuation exists.
Thus, we also use probability threshold-based evaluations, such as 80\%, where the editing success rate evaluates whether the loss $L_{\mathbf{x}_e, \mathbf{y}_{\text{ood}}}$ for an example falls below $\delta = - \text{log} (0.8)$, as indicated in the formula below.
The intuition behind this is that many other plausible alternative continuations may exist.\looseness=-1

\begin{equation}\label{equ:temporal_es}
    \vspace{-5pt}
    \text{OOD Gen.}  = \frac{1}{T}\sum\limits_{t=1}^{T} \mathbbm{1}\{(L_{\Theta_{T}}(\mathbf{x}_e, \mathbf{y}_{\text{ood}}) < \delta)\}.
\end{equation}

\begin{table}
    \centering
    \caption{\textbf{Temporal} OOD dataset example. Bolded text refers to the edit labels $\mathbf{y}_e$ and $\mathbf{y}_{\text{ood}}$.}
    \vspace{4pt}
    \begin{tabular}{p{1cm}p{12.3cm}}
        \toprule
        $\mathbf{x}_e, \mathbf{y}_e$ & \temporalx~\textbf{\temporaly} \\
        \cmidrule{2-2}
        $\mathbf{x}_{\text{loc}}$ & \temporallocx ~\temporallocy \\
        \cmidrule{2-2}
        $\mathbf{x}_e, \mathbf{y}_{\text{ood}}$ & \temporalx~\textbf{\temporalyy} \\
        \bottomrule
    \end{tabular}
    \label{tab:temporal-example}
    \vspace{-11pt}
\end{table}

\subsection{Descriptions of Compared Model Editors}\label{baselines}
\textbf{FT-L.} ~
All other layers of the LLMs remain frozen, and only a single MLP layer is fine-tuned through autoregressive loss \cite{rome}. 
Additionally, we impose an $\text{L}_{\infty}$ norm constraint to prevent the parameters from deviating too far from the pretrained distribution.

\textbf{FT-EWC.} ~
Elastic Weight Consolidation (EWC) has been demonstrated to mitigate catastrophic forgetting by updating weights using a Fisher information matrix, which is computed from past edits, multiplied by a scaling factor $\lambda$ \cite{ewc}.
Following \cite{grace}, we omit the constraints of the $\text{L}_{\infty}$ norm in this implementation.

\textbf{MEND.} ~
MEND \cite{mend} transforms the gradients obtained from standard fine-tuning using a hypernetwork that converts gradients decomposed into low rank (rank=1) into new gradients, which are then applied to the target layer for parameter updates.
During the training phase, a small auxiliary hypernetwork receives editing examples $(\mathbf{x}_e, \mathbf{y}_e)$, and $\mathbf{x}_{\text{loc}}$.
MEND's training loss comprises the standard autoregressive loss combined with the KL divergence loss of the model's output on $\mathbf{x}_{\text{loc}}$ before and after editing.
This hypernetwork plays a crucial role during the editing procedure.

\textbf{ROME.} ~
ROME \cite{rome} uses causal analysis to pinpoint knowledge within specific MLP layers and modifies the entire matrix through least squares approximation.
It operates under the strong assumption that the MLP is the primary module for storing knowledge \cite{geva-etal-2021-transformer}, and it injects a single piece of knowledge into the MLP at each iteration using a Lagrangian remainder.

\textbf{MEMIT.} ~
Similarly, based on the assumption that the FFN serves as a knowledge key-value store, MEMIT \cite{memit} manipulates parameters of specific layers directly through least squares approximation.
Unlike ROME, which updates a single layer, MEMIT is a multi-layer updating algorithm that supports simultaneous updates of hundreds or thousands of facts.
For sequential model editing tasks, MEMIT requires immediate on-the-fly repairs when the model makes errors, expressed as $f_{\Theta_{T}} = {\rm MEMIT}(f_{\Theta_{T-1}}, \mathbf{x}_T, \mathbf{y}_T),$ involving multiple operations on the original model.

\textbf{MEMIT-MASS.} ~
Unlike sequential editing, MEMIT supports modification of multiple knowledge fragments in a batch mode, named \textbf{MEMIT-MASS}.
Suppose we collect streaming errors as $(\mathcal{X}, \mathcal{Y}) = \{(\mathbf{x}_0, \mathbf{y}_0), (\mathbf{x}_1, \mathbf{y}_1), ..., (\mathbf{x}_T, \mathbf{y}_T)\}$ and inject them collectively into the MLP, it only involves a single editing operation on the original model as $f_{\Theta_{T}} = {\rm MEMIT}(f_{\Theta_{0}}, \mathcal{X}, \mathcal{Y}).$
Although this approach \textbf{loses the capability for on-the-fly repairs}, we still include this baseline in our experiments.

\textbf{DEFER.} ~
In GRACE, a reimplementation of SERAC \cite{serac} is utilized, denoted as DEFER.
For new inputs, DEFER includes a network $g$ (corresponding to the \textit{scope classifier} in SERAC) that predicts whether to: 1) trust the prediction of the LLMs, or 2) trust the prediction of the new model.
Here, the new model is configured as a single-layer linear network $o$ with a sigmoid activation function, corresponding to the \textit{counterfactual model} in SERAC.
During the editing process, $g$ and $o$ are fine-tuned jointly.

\textbf{GRACE.} ~
GRACE \cite{grace} utilizes a discrete KEY-VALUE codebook and maintains the codebook throughout the editing flow by adding, expanding, and splitting KEYs.
During the inference phase, it retrieves the nearest KEY and determines whether to replace the activation of the hidden layer output.

\subsection{Training Details and Hyperparameters} \label{app_sec: training_details}
Except for MEMIT-MASS, the batch size for all methods is consistently 1 in sequential editing scenarios.
All experiments are conducted using 3 NVIDIA A800 GPUs, with all tasks reproducible on a single A800.
Editing ZsRE takes approximately 4 hours, while Hallucination requires around 6 hours.
To ensure fair comparisons, unless otherwise specified (for some methods like MEND, ROME, and MEMIT, we follow the original literature by selecting the last few layers or using causal analysis to identify the target layers), the default target layers for editing on \texttt{LLaMA}, \texttt{Mistral}, and \texttt{GPT-J} are \texttt{model.layers[27].mlp.down\_proj.weight}, \texttt{model.layers[27].mlp.down\_proj.weight}, and \texttt{transformer.h[21].mlp.c\_fc}, respectively.

For FT-L, we utilize a reimplementation from ROME \footnote{\url{https://github.com/kmeng01/rome}}, employing the Adam \cite{adam} optimizer with consideration of learning rates at 1e-5, 1e-4, and 5e-4, and conducting gradient descents for 50 iterations, ultimately reporting the best results at a learning rate of 5e-4.

For FT-EWC, we follow the reimplementation in GRACE and its default settings, setting the learning rate at 1e-2, the $ \lambda_{\text{ewc}} $ penalty factor at 0.1, and the number of replay instances at 10.

\begin{figure}[!htbp]
    \begin{minipage}{1.0\linewidth}\centering
         % \vspace{-0.2cm}
    \centering
    \includegraphics[width=0.48\linewidth]{img/llama-causal-trace-no-attn-mlp.pdf}
    \includegraphics[width=0.48\linewidth]{img/mistral-causal-trace-no-attn-mlp.pdf}
    % \vspace{-13pt}%
    \caption{Mid-layer MLPs play a crucial mediating role in \texttt{LLaMA-2-7B} and \texttt{Mistral-7B}.}
    % \vspace{-13pt}%
    \label{fig:causal_analysis}
    \end{minipage}
\end{figure}
For the training phase of MEND, we adhere to the original paper, setting the learning rate at 1e-4, iterating 100K times, and employing early stopping at 30K, ultimately achieving an accuracy of 0.95 on the training set.
Notably, we target the last few MLP layers as per the original literature, such as \texttt{model.layers[i].mlp.down\_proj.weight}, \texttt{model.layers[i].mlp.gate\_proj.weight}, \texttt{model.layers[i].mlp.up\_proj.weight} in LLaMA, where $i \in [29, 30, 31]$.

For ROME and MEMIT, we follow the original literature on GPT-J using the default configurations, specifically the fifth layer and layers [3,4,5,6,7,8].
In LLaMA and Mistral, additional causal analysis is conducted to pinpoint the layers storing knowledge. As shown in Figure \ref{fig:causal_analysis}, an increasing trend in the Average Indirect Effect of the MLP is observed across layers [4,5,6,7,8], suggesting that the model recalls factual knowledge here and passes the matured token distribution via residual connections to the last MLP.
Thus, in LLaMA and Mistral, ROME edits the fifth layer, while MEMIT edits layers [4,5,6,7,8].

\begin{wraptable}{r}{0.35\textwidth}
  \centering
  \vspace{-0.4cm}
  \arrayrulecolor{black}
  \caption{WISE hyper-parameters during editing and merging.}
  \vspace{5pt}
  \resizebox{0.35\textwidth}{!}{
    \begin{tabular}{lr}
    \toprule
    Hyper-Parameters & Values \\
    \midrule
    Optimizer & \textsc{SGD} \\
    LR $\eta$ & $1.0$ \\
    Mask Ratio $\rho$ & $0.2$ \\
    $\alpha$ & $5.0$ \\
    $\beta$ & $20.0$ \\
    $\gamma$ & $10.0$ \\
    \midrule
    Merge Weights $\lambda$ & $0.5$ \\
    Knowledge shards $k$ & $2$ \\
    \bottomrule
    \vspace{-25pt}
    \end{tabular}}
  \label{tab:wise-hyper}%
\end{wraptable}
For DEFER, the original literature uses a learning rate of 1.0; however, we found it unfit for LLaMA and Mistral, with severe fluctuations in model loss.
Therefore, we experiment with learning rates of 7e-5, 7e-4, and 1e-3, and ultimately report using 7e-5 (optimal).

For GRACE, we strictly follow the original literature, setting the learning rate at 1.0, and using \texttt{replace\_last} to only replace the activation of the last token in autoregressive scenarios.
After observing failures in generalization, we adjust various $ \epsilon_{\text{init}} $ values and discuss this more in Appendix \ref{subsec: grace_generalization_analysis}.

For WISE, the hyperparameters for the QA and Hallucination tasks are identical.
We find that a learning rate of 1.0 with the SGD \cite{sgd} optimizer is a good approach for stable training.
The hyperparameters designed in the knowledge editing phase include the random masking probability $ \rho $ and the routing threshold guidance $ \alpha, \beta, \gamma $.
In the knowledge merging phase, hyperparameters include the number of merges $ k $ and the merging weights $ \lambda $ for each MLP (we discuss the impact of $ \rho $ and $ k $ in Section \ref{subsec: further_analysis}).
Theoretically, as the importance of knowledge in any MLP is considerable, we always average with $ \lambda = 1 / k $ across all experiments. These are shown in Table \ref{tab:wise-hyper}.

\subsection{Pseudo Code of WISE}
The pseudo-code of the WISE editing stage is in Algorithm~\ref{alg:wise_editing}, and the one of the WISE inference stage is Algorithm~\ref{alg:wise_inference}.

\begin{algorithm}[h]
\vspace{-3pt}
\caption{WISE Editing Stage}
\label{alg:wise_editing}
\vspace{0.1CM}
\textbf{Input}: The initial LLM model $f_{\Theta_0}$, the targeted FFN layer, the edit dataset $\mathcal{D}_{\text{edit}}$ whose length is $T$, the irrelevant dataset $\mathcal{D}_{\text{irr}}$, the subspace mask ratio $\rho$, the number of subspaces $k$, whether WISE-Retrieve.\\
\textbf{Output}: The final LLM model $f_{\Theta_T}$ after $T$ edits.
\begin{algorithmic}[1] %[1] enables line numbers
\STATE Generate $k$ random masks $\mathbf{M}_i, i \in [k]$ of ratio $\rho$; if WISE-Retrieve, copy the side memory several times;
\FOR{each edit $(\mathbf{x}_t, \mathbf{y}_t) \in \mathcal{D}_{\text{edit}}, t \in [T]$}
        \STATE Edit $(\mathbf{x}_t, \mathbf{y}_t)$ in the corresponding memory subspace by $L_{\text{edit}} = -\log P_{W_{v'}}(\mathbf{y}_t | \mathbf{x}_t) + L_a$;
        % \STATE Train the activation routing $\Delta_{\text{act}}$ according to Equation~\ref{equ:routing_loss} and update the threshold $\epsilon$;
        \STATE Update the activation threshold: $\epsilon = \min(\epsilon, \Delta_{\text{act}}(\mathbf{x}_t))$;
        \IF{All the $k$ subspaces of a side memory are full}
            \STATE Use Ties-Merge in Equation~\ref{equ:ties_update} to update the final side memory;
            \IF{WISE-Retrieve}
                \STATE Move to another copy of side memory $\mathbf{W}_{v'}$;
            \ENDIF
        \ELSE
            \IF{Current subspace $\mathbf{M}_i$ is full}
                \STATE Move to another subspace of side memory $\mathbf{M}_{i+1}$;
            \ENDIF
        \ENDIF
    \ENDFOR
\STATE \textbf{return} Obtain the final LLM model $f_{\Theta_T}$.
\end{algorithmic}
\end{algorithm}
\vspace{-9pt}


\begin{algorithm}[h]
\caption{WISE Inference Stage}
\label{alg:wise_inference}
\textbf{Input}: The edited LLM model $f_{\Theta_T}$, the activation threshold $\epsilon$, the test dataset $\mathcal{D}_{\text{test}}$, whether WISE-Retrieve.\\
\textbf{Output}: The model's output.
\begin{algorithmic}[1] %[1] enables line numbers
\FOR{each query $\mathbf{x}_i \in \mathcal{D}_{\text{test}}$}
        \IF{WISE-Retrieve}
            \STATE Get the value of activation $\Delta_{\text{act}} = \Vert \mathcal{A}(\mathbf{x}_i) \cdot (\mathbf{W}_{v'} - \mathbf{W}_v) \Vert_2$ for each side memory and select the one with the maximal value of $\Delta_{\text{act}}$;
            \ELSE 
                \STATE Get the value of activation $\Delta_{\text{act}} = \Vert \mathcal{A}(\mathbf{x}_i) \cdot (\mathbf{W}_{v'} - \mathbf{W}_v) \Vert_2$;
        \ENDIF
        \IF{$\Delta_{\text{act}} > \epsilon$}
            \STATE Use the side memory $\mathbf{W}_{v'}$ to generate the output as in Equation~\ref{equ:ffn_out};
        \ELSE
            \STATE Use the main memory $\mathbf{W}_{v}$ to generate the output as in Equation~\ref{equ:ffn_out}.
        \ENDIF
    \ENDFOR
\end{algorithmic}
\end{algorithm}


% \subsection{Additional Dataset}
% \paragraph{Measure editing OOD capabilities on Temporal} \label{temporal_metric}

% \begin{figure}[t]
%     \centering
%     \includegraphics[width=1.0\linewidth]{img/temporal_editing_example.pdf}
%     \caption{Given simple examples of model editing, we evaluate editing performance in more complex and natural settings (\textsc{OOD Prompt}). The Loc data is sampled from the pretraining corpus.}
%     \label{fig:temporal_example}
% \end{figure}


\section{More Experimental Results and Analyses} \label{app_sec:more_results}
% % \begin{table}[ht]
%     \centering
%     \setstretch{1.2}
%     \vspace{0.2cm}
%     \caption{\texttt{GPT-J-6B} Main editing results for ZsRE.}
%     \resizebox{0.9\linewidth}{!}{
%         \begin{tabular}{lccc|c|ccc|c|ccc|c|ccc|c}
%         \toprule
%         \multirow{3}{*}{\textbf{Method}}
        
%         %\cmidrule(lr){5-7}\cmidrule(lr){8-10} 
        
%          % \cmidrule(lr){2-7}
%          & \multicolumn{4}{c}{$\textbf{n = 1}$} & \multicolumn{4}{c}{$\textbf{n = 10}$} & \multicolumn{4}{c}{$\textbf{n = 100}$} & \multicolumn{4}{c}{$\textbf{n = 1000}$} \\
%          \cmidrule(lr){2-17}
%          & Reliability & Generality & Locality & Average & Reliability & Generality & Locality & Average & Reliability & Generality & Locality & Average & Reliability & Generality & Locality & Average\\
%          \midrule
%          FT-EWC & 1.000 & 0.9950 & 0.0020 & 0.6657 & 0.5726 & 0.5431 & 0.0064 & 0.3740 &0.5883 &0.5425 &0.0344 & 0.3884 &0.5172 &0.4546 &0.0079 & 0.3266 \\
%          MEND & 0.9723 & 0.9683 & 0.9788 & 0.9731 & 0.0117 & 0.0117 & 0.0075 & 0.0103 &0.0045 &0.0045 &0.000 & 0.0030 &0.0016 &0.0042 &0.000 & 0.0020 \\
%          ROME & 1.000 & 0.9638 & 0.8793 & 0.9477 & 0.4699 & 0.4470 & 0.1072 & 0.3414 &0.0214 &0.0173 &0.0262 & 0.0216 &0.0097 &0.0130 &0.0142 & 0.0123 \\
%          MEMIT-MASS  & 0.9875 & 0.8742 & 0.9877 & 0.9498 & 0.9763 & 0.8060 & 0.9796 & 0.9206 &0.7738 &0.7110 &0.9257 & 0.8035 &0.7568 &0.6995 &0.8601 & 0.7388 \\
%          DEFER  & 1.000 & 1.000 & 0.3886 & 0.7962 & 0.6618 & 0.5081 & 0.3599 & 0.5099 &0.0111 &0.0111 &0.9507 & 0.3243 &0.2627 &0.2524 &0.6961 & 0.4057 \\
%          GRACE & 1.000 & 0.0011 & 1.000 & 0.6670 & 0.9700 & 0.2700 & 1.000 & 0.7467 &0.98 &0.18 &1.00 & 0.7267 &0.99 &0.0669 &1.000 & 0.6856 \\
%          \textbf{WISE (ours)}  &1.000 &0.995 &1.000 & 0.9987 &0.9295 &0.9148 &1.000 & 0.9481 &0.8441 &0.7486 &1.00 & 0.8642 &0.7030 &0.6336 &1.000 & 0.7795 \\
%         \bottomrule 
%         \end{tabular}
%     }
%     \label{tab:gptj_zsre}
%     \vspace{-0.3cm}
% \end{table}

\begin{figure}[!t]
    \centering
    \includegraphics[width=0.245\textwidth]{img/gpt-j/gptj_zsre_Reliability.pdf}%
    \includegraphics[width=0.245\textwidth]{img/gpt-j/gptj_zsre_Generality.pdf}%
    \includegraphics[width=0.245\textwidth]{img/gpt-j/gptj_zsre_Locality.pdf}%
    \includegraphics[width=0.245\textwidth]{img/gpt-j/gptj_zsre_Average.pdf}%
    \vspace{-3pt}%
    \caption{GPT-J-6B, ZsRE, continual editing.}
    \label{fig:gpt-j_zsre_fig}
\end{figure}

% data = {
%     "reliability": {
%         "FT-EWC": [1.000, 0.5726, 0.5883, 0.5172],
%         "MEND": [0.9723, 0.0117, 0.0045, 0.0030],
%         "ROME": [1.000, 0.4699, 0.0214, 0.0216],
%         "MEMIT-MASS": [0.9875, 0.9763, 0.7738, 0.8035],
%         "DEFER": [1.000, 0.6618, 0.0111, 0.3243],
%         "GRACE": [1.000, 0.9700, 0.9800, 0.9900],
%         "WISE (ours)": [1.000, 0.9295, 0.8441, 0.7030]
%     },
%     "generality": {
%         "FT-EWC": [0.9950, 0.5431, 0.5425, 0.4546],
%         "MEND": [0.9683, 0.0117, 0.0045, 0.0016],
%         "ROME": [0.9638, 0.4470, 0.0173, 0.0097],
%         "MEMIT-MASS": [0.8742, 0.8060, 0.7110, 0.7568],
%         "DEFER": [1.000, 0.5081, 0.0111, 0.2627],
%         "GRACE": [0.0011, 0.2700, 0.1800, 0.0669],
%         "WISE (ours)": [0.995, 0.9148, 0.7486, 0.6336]
%     },
%     "locality": {
%         "FT-EWC": [0.0020, 0.0064, 0.0344, 0.0079],
%         "MEND": [0.9788, 0.0075, 0.0000, 0.0042],
%         "ROME": [0.8793, 0.1072, 0.0262, 0.0130],
%         "MEMIT-MASS": [0.9877, 0.9796, 0.9257, 0.6995],
%         "DEFER": [0.3886, 0.3599, 0.9507, 0.2524],
%         "GRACE": [1.000, 1.000, 1.000, 1.000],
%         "WISE (ours)": [1.000, 1.000, 1.000, 1.000]
%     },
%     "average": {
%         "FT-EWC": [0.6657, 0.3740, 0.3884, 0.3266],
%         "MEND": [0.9731, 0.0103, 0.0030, 0.0020],
%         "ROME": [0.9477, 0.3414, 0.0216, 0.0123],
%         "MEMIT-MASS": [0.9498, 0.9206, 0.8601, 0.7388],
%         "DEFER": [0.7962, 0.5099, 0.3243, 0.4057],
%         "GRACE": [0.6670, 0.7467, 0.7267, 0.6856],
%         "WISE (ours)": [0.9987, 0.9481, 0.8642, 0.7795]
%     }
% }


% activation scatter
% size scaling
% grace generality (t-SNE) [meaning trajectory]
% Hparams: Mask_Ratio, Merge_Freq
% Layer粒度，为什么选择中后层
% TRR 等价于 locality
% RAG 为什么泛化性差，ICL不efficiency。
% 和SERAC的具体区别
\begin{figure}[!t]
    \centering
    \includegraphics[width=0.245\textwidth]{img/gpt-j/gptj_zsre_Reliability.pdf}%
    \includegraphics[width=0.245\textwidth]{img/gpt-j/gptj_zsre_Generality.pdf}%
    \includegraphics[width=0.245\textwidth]{img/gpt-j/gptj_zsre_Locality.pdf}%
    \includegraphics[width=0.245\textwidth]{img/gpt-j/gptj_zsre_Average.pdf}%
    \vspace{-3pt}%
    \caption{GPT-J-6B, ZsRE, continual editing.}
    \label{fig:gpt-j_zsre_fig}
    \vspace{-3pt}
\end{figure}

\subsection{On the Pitfall of GRACE: Generalization Collapses in Decoder-only LLMs} \label{subsec: grace_generalization_analysis}

\begin{figure}[t]
    \begin{minipage}{0.99\textwidth}
        \centering
        \includegraphics[width=\linewidth]{img/radii_1.pdf}
        % \caption{Block 27. $\epsilon_\text{init} = 1.0$}
        \label{fig:radii_1.0}
        \vspace{-0.78cm}
    \end{minipage}
    \begin{minipage}{0.99\textwidth}
        \centering
        \includegraphics[width=\linewidth]{img/radii_3.pdf}
        % \caption{Block 27. $\epsilon_\text{init} = 3.0$}
        \label{fig:radii_3.0}
        \vspace{-0.78cm}
    \end{minipage}
    \begin{minipage}{0.99\textwidth}
        \centering
        \includegraphics[width=\linewidth]{img/radii_10.pdf}
        % \caption{Block 27. $\epsilon_\text{init} = 10.0$}
        \label{fig:radii_10.0}
        \vspace{-0.78cm}
    \end{minipage}
    \begin{minipage}{0.99\textwidth}
        \centering
        \includegraphics[width=\linewidth]{img/radii_500.pdf}
        % \caption{Block 27. $\epsilon_\text{init} = 500.0$}
        \label{fig:radii_500.0}
    \end{minipage}
    \vspace{-0.5cm}
    \caption{Investigation on the query $\mathbf{x}$ and its distance to the nearest Key $k$, as well as the \textit{deferral radius} $\epsilon$ of that Key. Red and Blue respectively represent the paraphrase query $\mathbf{x}_{e'}$ and the unrelated query $\mathbf{x}_{\text{loc}}$, with the hatch representing the radius of the nearest Key. We observe that when conflicts occur (hit the codebook Key but with different Edit Target $\mathbf{y}_e$), the \textit{deferral radius} $\epsilon$  decays exponentially. This results in GRACE being unable to encompass the paraphrase $\mathbf{x}_{e'}$ and maintain high locality, regardless of how $\epsilon_{\text{init}}$ is adjusted. ZsRE, \texttt{LLaMA-2-7B}.}\label{fig:grace_failure}
\end{figure}

Here, we discuss why GRACE exhibits poor generalization when editing decoder-only LMs.

As shown in Figure \ref{fig:grace_failure}, we continuously edit 15 samples $(\mathbf{x}_e, \mathbf{y}_e)$ using GRACE and observe the nearest codebook \textit{Key} for their paraphrases $\mathbf{x}_{e'}$ and unrelated queries $\mathbf{x}_{\text{loc}}$, as well as the governed \textit{Deferral radii} $\epsilon$ of those \textit{Keys}.
When overlapping \textit{Keys} exist, GRACE reduces the \textit{Deferral radii} to split this \textit{Keys} and then adds a new codebook entry, resulting in exponentially decaying of radii $\epsilon$ during the editing process. Though $\epsilon$ is initialized from a high $\epsilon_{\text{init}}$, it will be small and ineffective after continuous edits. 
From Figure \ref{fig:grace_failure}, we observe that GRACE is more likely to have a conservative strategy that sets smaller Deferral radii during editing. Smaller Deferral radii will cause $\mathbf{x}_{e'}$ to fail to hit the codebook (the distance to the nearest \textit{Key} is farther than its \textit{Deferral radii}) but let $\mathbf{x}_{\text{loc}}$ successfully far away from the radii, resulting low generalization and high locality. Also, we observe that the Deferral radii method is not effective under any $\epsilon_{\text{init}}$; for all tested $\epsilon_{\text{init}}$ values of 1.0, 3.0, 10.0, and 500.0, they all have low generalization and high locality. 

This suggests that in autoregressive LMs, the distribution of the last token cannot effectively represent semantics; whereas in encoder-only and encoder-decoder architectures, capturing semantic information through vector representation has been extensively studied \cite{bert, sentencebert, simcse}. This is consistent with the degree of generalization shown by GRACE when anchoring the T5 \cite{T5} Encoder layer. 
Some related works \cite{meaning_trajectory} also indicate that in autoregressive models, semantic similarity measures based on averages of output tokens underperform, recommending the use of score distributions over text continuations to represent semantic distances.

\subsection{Impact of Knowledge Merging Strategies for WISE} \label{app_sec:varying_merge}
\begin{wraptable}{r}{0.3\textwidth}
    \centering
    % \vspace{-1.0cm}
    \caption{\small Varying Merging Strategy. ZsRE. \texttt{LLaMA-2-7B}.}
    \resizebox{1\linewidth}{!}{
    \begin{tabular}{l|cccc}
    \toprule
    \textbf{Methods} & Rel. & Gen. & Loc. & Avg.\\
    \midrule
    \textit{Linear}  & .63 & .61 & .93 & .72\\
    \textit{Slerp}  & .62 & .64 & .91  & .72\\
    \textit{Dare}  & .68 & .63 & .92  & .74\\
    \textit{Dare\_Ties}  & .67 & .63 & .83 & .71\\
    \rowcolor[gray]{0.94} 
    \textit{Ties}  & \textbf{.85} & \textbf{.81} & \underline{.94}  & \textbf{.87}\\
    \rowcolor[gray]{0.94} 
    \textit{Sign}  & \underline{.80} & \underline{.76} & \textbf{.97} & \underline{.84}\\
    \bottomrule
    \end{tabular}}
    \vspace{-10pt}
    \label{tab:knowledge-merging-strategy}
\end{wraptable}

Here, we conduct a more in-depth study of the knowledge merging strategies for WISE, exploring various merging approaches including ($i$)~\underline{\textit{Linear}}, which uses a simple weighted average;
($ii$)~\underline{\textit{Slerp}}, which spherically interpolates the parameters of two models;
($iii$)~\underline{\textit{Ties}}, a component used in the main experiments of this paper that resolves merging disturbances through TRIM ELECT SIGN;
($iv$)~\underline{\textit{Dare}}: which follows a Bernoulli distribution to delete redundant parameters and rescale the remaining ones;
($v$)~\underline{\textit{Dare\_Ties}}, which combines dare and the sign consensus algorithm of TIES;
and ($vi$)~\underline{\textit{Sign}}, an ablation component of Ties that addresses directional conflicts—all utilizing the official implementation from MergeKit \cite{mergekit} \footnote{\url{https://github.com/arcee-ai/mergekit}}.
We randomly sample 100 edits from ZsRE, retaining a fine-tuned MLP every 50 edits (merging 2 MLPs).
As shown in Table \ref{tab:knowledge-merging-strategy}, we observe that ignoring the direction of parameter updates (Linear, Slerp, Dare) leads to a significant decline in editing performance, underscoring the importance of addressing knowledge conflicts in overlapping parameters.
The success of \textit{Sign} also reaffirms this point.
Meanwhile, the randomly masked knowledge shards exhibit a non-redundancy, indivisible nature.
This is demonstrated by the significantly weaker performance of \textit{Dare\_Ties} compared to \textit{Ties}/\textit{Sign}, indicating that removing parameter updates can lead to the loss of edited knowledge or even potential "anchors".

\subsection{Analysis of Retrieving Top-1 Activation} \label{app_sec: retrieve_analysis}
\begin{figure}[!htbp]
    \centering
    \begin{subfigure}[t]{0.35\linewidth}
        \centering
        \includegraphics[width=\linewidth]{img/memo_act_es.pdf}
        \caption{\hspace{5pt}Average of Rel. and Gen.}
        \label{fig:memo_act_es}		
    \end{subfigure}
    \hspace{5pt}
    \begin{subfigure}[t]{0.358\linewidth}
        \centering
        \includegraphics[width=\linewidth]{img/memo_act_prec.pdf}
        \caption{\hspace{3pt}Retrieval Acc. by Top-1 Activation}
        \label{fig:memo_act_prec}	
    \end{subfigure}
    % \centering
    % \includegraphics[width=0.7\linewidth]{img/memo_act.pdf}
    % \vspace{-13pt}%
    \caption{Comparing editing results of WISE-\texttt{\{Retrieve, $\text{Retrieve}_\text{oracle}$, Retrieve w. $\mathrm{L}_{\text{memo}}$\}} when varying $T$. (a) shows the simple average of Rel. and Gen. (ES.), while 
    (b) shows retrieval accuracy, i.e., whether the Top-1 Activation routes to the correct MLP (prec@1). \texttt{X-axis}: Num edits. ZsRE. \texttt{LlaMA-2-7B}.}
    % \vspace{-13pt}%
    \label{fig:memo_act}
\end{figure}

WISE-\texttt{Retrieve} retains each knowledge-sharding memory and retrieves through Top-1 Activation.
However, as shown in Table \ref{tab:zsre_scaling2_3K} and Figure \ref{fig:memo_act_prec}, the retrieval accuracy still has significant room for improvement; specifically, when $T$ reaches 3K, the accuracy of routing to the correct MLP drops to around 60\%, indicating the specificity between side memories is insufficient.
One possible reason is that when sampling the edits from a single dataset (ZsRE), the editing instances $(\mathbf{x}_e, \mathbf{y}_e)$ all belong to the same domain.
This leads to some very similar instances being captured by multiple expert side memories (resulting in high activations for all side memories), introducing more retrieval failures.

Therefore, to improve the specificity of side memory and reduce the probability of routing errors, we attempt to add a new constraint $L_{\text{memo}}$ to Equation \ref{equ:routing_loss}.
For knowledge-sharding memory $\mathbf{W}_i$, we randomly replay instances $(\mathbf{x}_{\mathrm{m}}, \mathbf{y}_{\mathrm{m}})$ from the edit set $\mathcal{D}_{\mathbf{W}_j}$ of past shard $\mathbf{W}_{j, \, j \in [0, i-1]}$, ensuring that $\mathbf{W}_i$ remains inactive for $\mathbf{x}_{\mathrm{m}}$:

$$L_a' = L_a + \underbrace{\max(0, \Delta_{\text{act}}(\mathbf{x}_{\mathrm{m}}) - \alpha)}_{L_\text{memo}}, \quad \text{s.t.}\; \mathbf{x}_{\mathrm{m}} \in \mathcal{D}_{\mathbf{W}_j}.$$

As shown in Figure \ref{fig:memo_act_prec}, this replay behavior increases the specificity between side memories, maintaining nearly \textbf{88\%} retrieval accuracy at $T=3K$. 
Figure \ref{fig:memo_act_es} also shows that WISE-\texttt{Retrieve} \textit{w.} $L_{\text{memo}}$ improves Edit Success (ES.) by \textbf{8.39\%} compared to WISE-\texttt{Retrieve}, providing a promising direction for future work.
With finer-grained activation management, we might be able to bridge the performance gap between \texttt{Retrieve} and \texttt{Oracle}.

\subsection{Case Study}
\newcommand{\greenyes}{\textcolor{green}{\ding{51}}}
\newcommand{\bluehalf}{\textcolor{cyan}{\ding{52}\rotatebox[origin=c]{-9.2}{\kern-0.7em\ding{55}}}}
\newcommand{\redno}{\textcolor{red}{\ding{55}}}

\begingroup
\setlength{\tabcolsep}{5pt} % Default value: 6pt
\begin{table}[!htbp]
    % \color{blue}
    \small
    \centering
    \setstretch{1.2}
    \vspace{-5mm}
    \caption{\textbf{Failure cases of using WISE} to edit \texttt{LLaMA-2-7B}. \bluehalf represents errors in part of the tokens, \redno represents complete output errors (i.e., factual failures), and \greenyes indicates the expected exact match.}
    \vspace{3pt}
    % \begin{tabular}{l@{\hspace{2pt}}ccc}
    \begin{tabular}{l@{\hspace{2pt}}>{\raggedright\arraybackslash}p{0.48\textwidth}>{\raggedright\arraybackslash}p{0.18\textwidth}>{\raggedright\arraybackslash}p{0.19\textwidth}}
        \toprule
        & \normalsize Prompt & \normalsize Edit Target & \normalsize Post-Edit Output \\
        \midrule
        $i$\textit{a}) & By which person Lahti Town Hall has been designed? & Aki Kaurismäki & Wime Kaurismäki  \bluehalf \\
        $i$\textit{b}) & \textit{Which is the architect of Lahti Town Hall?} & - & Wime Kaurismäki  \bluehalf \\
        $i$\textit{c}) & Which corporation was USS Leedstown (APA-56) created by? & Lockheed Shipbuilding & Leez Shipbuilding \bluehalf \\
        $i$\textit{d}) & \textit{Which company manufactures the USS Leedstown (APA-56)?} & - & Leez Shipbuilding \bluehalf \\
        \midrule
        $ii$\textit{a}) & Which language is Garowe Principles written in? & Persian & Dutchian  \redno \\
        $ii$\textit{b}) & \textit{In what language does the monthly football magazine Garowe Principles report?} & - & Somian  \redno \\
        $ii$\textit{c}) & What year was the service entry date for Panzer 58? & 1957 & 1953 \redno \\
        $ii$\textit{d}) & \textit{What was the year Panzer 58 was commissioned?} & - & 1953 \redno \\
        \midrule
        $iii$\textit{a}) & What was Gemma Bosini's range? & mezzo-srano & Wzo-srano  \redno \\
        $iii$\textit{b}) & \textit{The kind of voice of Gemma Bosini is what?} & - & mezzo-srano  \greenyes \\
        \midrule
        $iv$\textit{a}) & In which state is Qaleh Lan located? & Golestan Province & Golestan Province  \greenyes \\
        $iv$\textit{b}) & \textit{What state is Qaleh Lan in?} & - & Lestan Province  \redno \\
        $iv$\textit{c}) & In which language Garowe Principles monthly football magazine reporting? & American English & American English \greenyes \\
        $iv$\textit{d}) & \textit{What language are Garowe Principles written in?} & - & English English \redno \\
        \bottomrule
    \end{tabular}
    \label{tab:wise_fail_examples}
\end{table}
\endgroup

% \begingroup
% \setlength{\tabcolsep}{4pt} % Default value: 6pt
% \begin{table}
%     % \color{blue}
%     \small
%     \centering
%     \vspace{5mm}
%     \begin{tabular}{l@{\hspace{2pt}}>{\raggedright\arraybackslash}p{0.37\textwidth}>{\raggedright\arraybackslash}p{0.21\textwidth}>{\raggedright\arraybackslash}p{0.10\textwidth}>{\raggedright\arraybackslash}p{0.20\textwidth}}
%         \toprule
%         & \normalsize Input & \normalsize Pre-Edit Output & \normalsize Edit Target & \normalsize Post-Edit Output \\
%         \midrule
%         1a: & \textbf{Who is the president of the USA?} & Donald Trump \redno & \textbf{Joe Biden} & Joe Biden \greenyes \\
%         1b: & Who is the US president? & David Rice Atchison \redno & - & Joe Biden \greenyes \\
%         1c: & Who is the president of France? & Emmanuel Macron \greenyes & - & Emmanuel Macron \greenyes \\
%         \midrule
%         2a: & \textbf{Who designed the Burj Khalifa?} & British architect Herbert Baker \redno & \textbf{Skidmore, Owings \& Merrill} & Skidmore, Owings \& Merrill \greenyes \\
%         2b: & Who designed the Eiffel Tower? & Alexandre Gustave Eiffel \greenyes & - & Alexandre Gustave Eiffel \greenyes \\
%         2c: & Who designed the Empire State Building? & Shreve, Lamb and Harmon \greenyes & - & Shreve, Lamb and Harmon \greenyes \\
%         2d: & Who designed the Sydney Opera House? & Jrn Oberg Utzon \greenyes & - & Jrn Oberg Utzon$^*$ \greenyes \\
%         2e: & What firm was behind the design for the Burj Khalifa? & McKim, Mead \& White \redno & - & Skidmore, Owings \& Merrill \greenyes \\
%         2f: & What firm did the Burj Khalifa? & Jumeirah Group \redno & - & Jumeirah Group \redno \\
%         \midrule
%         3a: & \textbf{What car company makes the Astra?} & Mahindra \redno & \textbf{Opel} & Opel \greenyes \\
%         3b: & What car company makes the Mustang? & Ford \greenyes & - & Ford \greenyes \\
%         3c: & What car company makes the Model S? & Tesla Motors \greenyes & - & Tesla \greenyes \\
%         3d: & What car company makes the Wrangler? & Jeep \greenyes & - & Jeep \greenyes \\
%         3e: & What car company makes the F-150? & Ford \greenyes & - & Opel \redno \\
%         3f: & What car company makes the Golf? & Volkswagen AG \greenyes & - & Opel \redno \\
%         \midrule
%         4a: & \textbf{What artist recorded Thriller?} & Madonna \redno & \textbf{Michael Jackson} & Michael Jackson \greenyes \\
%         4b: & What artist recorded Dark Side of the Moon? & Pink Floyd \greenyes & - & Pink Floyd \greenyes \\
%         4c: & What artist recorded Bridge over Troubled Water? & Simon \& Garfunkel \greenyes & - & Simon \& Garfunkel \greenyes \\
%         4d: & What artist recorded Hotel California? & Don Henley \yellowmaybe & - & Don Henley \yellowmaybe \\
%         4e: & What band recorded Back in Black? & AC/DC \greenyes & - & Michael Jackson \redno \\
%         \bottomrule
%     \end{tabular}
%     \caption{\textbf{Additional examples of using MEND} to edit a 770M parameter. Example 2e shows correct generalization behavior; 2f shows an instance of \textbf{undergeneralization}; examples 3e, 3f, and 4e show instances of \textbf{overgeneralization}.}
%     \label{tab:examples-appendix}
% \end{table}
% \endgroup

%元组失败的；单个token的；难以泛化的，惊讶的可以泛化的；未能保持原分布的activation为啥寄了？

In Table \ref{tab:wise_fail_examples}, we present bad cases of using WISE to edit the \texttt{LLaMA-2-7B} on the ZsRE dataset and mitigating these failures is critical for future work in model editing.
We observe that in $i)$ errors occur only in part of the tokens, and these errors constitute a large proportion of the bad cases, indicating that the edits have not been sufficiently fitted.
$ii)$ displays cases where the entire output is incorrect, and factual failures indicate difficulties in retaining memory of parameters for some rare entities (such as Persian $iia, iib$).
$iv)$ presents cases of generalization failure, for example in $ivd)$, where the model answered ``English'' but did not fully follow the ground truth, indicating significant room for improvement in the accuracy of generalized edits.
Meanwhile, in $iii)$ we surprisingly find that even when WISE errs on the Edit Prompt, it can correctly answer its paraphrase $iii b)$ \textit{``The kind of voice of Gemma Bosini is what?''}.
This indicates that WISE can handle contextual information correctly in some cases but falls short in specific editing instructions, suggesting that optimizing editing instructions (modifying the editing context) may be a direction for improvement.


\subsection{Importance of \textit{Knowledge Anchor} When Merging Models}\label{app_subsec:knowledge_anchor}

\begin{wraptable}{l}{0.4\textwidth}
    \definecolor{gray}{rgb}{0.9373,0.9373,0.9373}
    \newcolumntype{a}{>{\columncolor{gray}}c}
    \centering
    \vspace{-0.5cm}
    \caption{\small Analysis of Merging \textit{w.o.} and \textit{w.} "knowledge anchor" (KA). $T=1000$. ZsRE. \texttt{LLaMA-2-7B}.}
    \resizebox{1.0\linewidth}{!}{
        \begin{tabular}{lcccc|aaaa}
        \toprule
        \multirow{2}{*}{\textbf{$\rho/k$}}
        
         %\cmidrule(lr){5-7}\cmidrule(lr){8-10} 
         % \cmidrule(lr){2-7}
         & \multicolumn{4}{c}{\textit{w.o.} KA} & \multicolumn{4}{c}{\textit{w.} KA} \\
         \cmidrule(lr){2-9}
         & \multicolumn{4}{c|}{\makecell{Rel.  ~Gen.  ~Loc. ~Avg.}} & \multicolumn{4}{c}{\makecell{Rel.  ~Gen.  ~Loc. ~Avg.}} \\
         \midrule
        2/0.30 & 0.76 & 0.72 & 1.00  &0.83 & \textbf{0.79} & \textbf{0.73} & 1.00 & \textbf{0.84}  \\
        2/0.50 &0.74 &\textbf{0.73} &1.00 & 0.82 &\textbf{0.77} &0.72 &1.00 &\textbf{0.83} \\
        3/0.33 &0.72 &0.68 &1.00 &0.80 &\textbf{0.75} &\textbf{0.71} &1.00 &\textbf{0.82} \\
        5/0.20 &0.64 &0.61 &1.00 &0.75 &\textbf{0.73} &\textbf{0.68} &1.00 &\textbf{0.80} \\
        \bottomrule 
        \end{tabular}
    }
    \label{tab:KA_ablation}
    \vspace{-0.4cm}
\end{wraptable}

Here, we discuss the effects of independent (ensured by non-overlapping masks) vs partially overlapping parameters within MLP subspaces on editing performance, as shown in Table \ref{tab:KA_ablation}.
It is observable that, despite varying mask ratios \(\rho\) and the number of subspaces \(k\), partial overlap (w. KA) consistently outperforms independent configurations (w.o. KA) in terms of Reliability (Rel.) and Generalization (Gen.). 
For example, at \(\rho/k\) of 5/0.20, there is a relative improvement of 9\% and 7\% respectively.
This demonstrates that the overlapping regions contribute as ``anchors'' for knowledge fusion, facilitating information transfer across different subspaces.
Moreover, the shared parameters provide a natural regularization \cite{subspace_reg} mechanism, helping synchronize model behavior across different subspaces.

\subsection{Ablation Study of Random Prefix Token}
\label{sec:random_token_ablation}
\begin{figure}[tbh!]
    \vspace{-0.3cm}
    \definecolor{figgreen}{rgb}{0.547,0.695,0.434}
    \definecolor{figblue}{rgb}{0.145, 0.451, 0.73}
    \definecolor{figyellow}{rgb}{1.000, 0.547, 0}
    \centering
    \includegraphics[width=0.8\linewidth]{img/pt_ablation.pdf}
    \vspace{-0.2cm}
    \caption{\textbf{Ablation studies on Random Prefix Token (PT) of WISE.} Light/Dark colors indicate the Editing Sucess w.o./w. PT addition. ZsRE. \texttt{LlaMA-2-7B}}
    \label{fig:random_token_ablation}
    \vspace{-11pt}%
\end{figure}

As described in Section \ref{subsec:exp_setting}, we employ random prefix token augmentation to enable the editing knowledge to cope with various contexts. That is, for a single $\mathbf{x}_e$, it expands into $
(\text{prefix}_{i}, \mathbf{x}_e)$. The prefix is derived from tokens that are randomly generated by the original LM $f_{\Theta}$, serving as an economical data augmentation method. We observe that the editing success rate is compromised (Figure \ref{fig:random_token_ablation}). Specifically, for instance, at T=1000, Rel. and Gen. decreased by 0.15 and 0.17, respectively. By utilizing randomly generated prefix tokens, the model is able to learn a broader range of linguistic features, thereby exhibiting greater robustness in practical applications. We believe that access to the "data generator" can deepen the model's memory of editing samples.




\subsection{Parameter Efficiency} \label{subsec: parameter_efficiency}
\begin{wrapfigure}{r}{0.4\columnwidth}
    \centering
    \vspace{-23pt}
    \includegraphics[width=0.4\textwidth]{img/gpu_consumption.pdf}
    \vspace{-17pt}%
    \caption{Computational costs.}
    \label{fig:gpu_consum}
    \vspace{-13pt}%
\end{wrapfigure}
The key to lifelong model editing is maintaining constant or slowly increasing computational costs as the number of edits expands.
Here, we provide a quantitative analysis using \texttt{LLaMA-2-7B} as an example.
Suppose we select \texttt{model.layers[27].mlp.down\_proj.weight} as side memory. In that case, the theoretically added parameters are \(11008 \times 4096 \times 4 = 0.18\) GB, which accounts for 0.64\% of the original LLaMA's \(7B \times 4 = 28\) GB (ignoring the VRAM required for input activations).
As shown in Figure \ref{fig:gpu_consum}, in practice, WISE-Merge increases VRAM by 4\% compared to the original \texttt{LLaMA} and remains constant over time.
WISE-Retrieve, instead of merging, uses retrieval routing, meaning the computational cost increases over time, but this increase is gradual and can easily handle thousands or tens of thousands of inputs.
Additionally, if we partially merge side MLPs (combining WISE-Retrieve and WISE-Merge), we can further reduce the computational demands of WISE-Retrieve.
\section{Proof of Theorem~\ref{theorem:subspace_overlap}}\label{app_sec:proof}

\newcommand{\qedsymbol}{\hfill $\square$}

% \begin{theorem}\label{app_theorem:subspace_overlap}
%     \textbf{Subspace Overlap.} 
%     Generate $k$ memory subspaces $\mathbf{W}_{v'}^{i}, i \in [k]$ by random mask with 1's ratio $\rho$, so each memory has $\rho\cdot|\mathbf{W}_{v'}|$ active trained parameters. For any two subspaces $\mathbf{W}_{v'}^{i}$ and $\mathbf{W}_{v'}^{j}$ $i\neq j ;i,j \in [k]$, there are $\rho^2\cdot|\mathbf{W}_{v'}|$ active parameters that are overlapped. For all $k$ subspaces, there are $\rho^k\cdot|\mathbf{W}_{v'}|$ overlapped active parameters.
% \end{theorem}

% \textit{Proof:} Different random masks are sampled independently, so we treat the sampling of the masks as stochastic independent events. Therefore, we have the following induction.

% \begin{enumerate}
%     \item The probability that one parameter is sampled each time is $\rho$.
%     \item If we sample twice because the sampling is independent, the probability that one parameter is sampled twice is $\rho^2$.
%     \item $\dots$
%     \item If we sample $k$ times, the probability that one parameter is sampled $k$ times is $\rho^k$.
% \end{enumerate}

% The number of parameters is $|\mathbf{W}_{v'}|$. Thus, the number of parameters overlapped in two random masks is $\rho^2\cdot|\mathbf{W}_{v'}|$ and the number of parameters overlapped in $k$ random masks is $\rho^k\cdot|\mathbf{W}_{v'}|$.

% \qedsymbol

\begin{theorem}\label{app_theorem:subspace_overlap}
    \textbf{Subspace Overlap.} 
    Generate $k$ memory subspaces $\mathbf{W}_{v'}^{i}, i \in [k]$ by random mask with 1's ratio $\rho$, so each memory has $\rho\cdot|\mathbf{W}_{v'}|$ active trained parameters. For any two subspaces $\mathbf{W}_{v'}^{i}$ and $\mathbf{W}_{v'}^{j}$ $i\neq j ;i,j \in [k]$, there are $\rho^2\cdot|\mathbf{W}_{v'}|$ active parameters that are overlapped. For all $k$ subspaces, there are $\rho^k\cdot|\mathbf{W}_{v'}|$ overlapped active parameters.
\end{theorem}

\textit{Proof:} We aim to prove the Subspace Overlap theorem by induction.

Let $\mathbf{W}_{v'}^{i}$ represent the $i$-th memory subspace generated by a random mask with a sparsity ratio of $\rho$, where $i \in [k]$. Each memory subspace $\mathbf{W}_{v'}^{i}$ contains $\rho\cdot|\mathbf{W}_{v'}|$ active trained parameters.

We start by considering the case of two memory subspaces, $\mathbf{W}_{v'}^{i}$ and $\mathbf{W}_{v'}^{j}$, where $i \neq j$ and $i, j \in [k]$.

Let $P(\text{parameter sampled}) = \rho$ be the probability that a parameter is sampled in one mask generation event.

\begin{enumerate}
    \item For a single mask generation, the probability that a specific parameter is sampled is $\rho$. We denote this probability as $P(\text{sampled}) = \rho$.
    \item Considering two independent mask generation events, the probability that the same parameter is sampled in both masks is the product of their individual probabilities, i.e., $\rho^2$. This is derived from the independence of the events. Mathematically:
    \[ P(\text{sampled in both masks}) = P(\text{sampled}) \times P(\text{sampled}) = \rho \times \rho = \rho^2. \]
    \item Extending this logic, for $k$ independent mask generation events, the probability that a specific parameter is sampled in all $k$ masks is $\rho^k$. Mathematically:
    \[ P(\text{sampled in all } k \text{ masks}) = \underbrace{P(\text{sampled}) \times P(\text{sampled}) \times \cdots \times P(\text{sampled})}_{k \text{ times}} = \rho^k. \]
\end{enumerate}

Now, let's calculate the number of parameters overlapped in two random masks:

The total number of parameters in $\mathbf{W}_{v'}$ is $|\mathbf{W}_{v'}|$.

Thus, the number of parameters overlapped in two random masks, $\mathbf{W}_{v'}^{i}$ and $\mathbf{W}_{v'}^{j}$, is $\rho^2\cdot|\mathbf{W}_{v'}|$.

Extending this to $k$ random masks, the number of parameters overlapped in all $k$ masks is $\rho^k\cdot|\mathbf{W}_{v'}|$.

This concludes the proof.

\qedsymbol


\section{Detailed Related Works}\label{app_sec:more_related_works}
\paragraph{Memory and Knowledge Injection of LLMs} 
The memories of LLMs can be divided into long-term (episodic) memory and working memory (short-term)~\cite{llm_controllable_working_memory,memoria,spalm}. Long-term memory refers to the knowledge stored in the model's parameters, which can be updated by (re)pretraining~\cite{gpt-1}, finetuning~\cite{lora}, and model editing~\cite{knowledge_editing_survey_1}. Working memory is stored in sustained activations/representations of neurons, which will be awakened during inference time~\cite{llm_controllable_working_memory}. In-context learning (ICL) is a kind of working memory~\cite{chain_of_thoughts}, also along with retrieval-based editing methods like GRACE~\cite{grace}. 
How to reinforce memory and inject/update knowledge for LLMs is a fundamental question~\cite{memoryllm,ft_vs_rag,fewshotft_vs_icl}. ICL or finetuning? Different works show different conclusions. In \cite{fewshotft_vs_icl}, the authors find that few-shot finetuning is more generalizable than ICL, especially for out-of-distribution data. In \cite{ft_vs_rag}, the authors contrast finetuning with retrieval-augmented generation (RAG) in terms of knowledge injection and find that RAG is better in most cases, and combining both will produce the best results. 
However, finetuning and pretraining are computation-expensive~\cite{llama-2,grace} and usually suffer from catastrophic forgetting~\cite{catastrophic_forgetting_llm} and overfitting~\cite{memorization_overfitting}. For ICL and RAG, the working memory is sometimes not controllable, the model may not follow the information of the contexts~\cite{llm_controllable_working_memory}, and the context window is limited~\cite{landmark_attention,infini_attention}, and there are works addressing these issues by training controllable ICL~\cite{llm_controllable_working_memory}, long-context ~\cite{landmark_attention,infini_attention}, and recurrent memory architecture design~\cite{memoryllm}. 
SPALM is proposed to add language models with storage modules that resemble both working and long-term memories~\cite{spalm}. 
% Also, model editing is an emerging technology that aims to enable data-efficient, fast, and non-destructive knowledge updates to LLMs~\cite{knowledge_editing_survey_1,knowledge_editing_survey_2}.

\paragraph{Model Editing of LLMs} 
Model editing can be summarized as the following lines of research. 
\textbf{\textit{Constrained finetuning:}} Preliminary model editing uses constrained finetuning to update parameters based on new examples~\cite{constrained_ft_me_1,constrained_ft_me_2}. 
\textbf{\textit{Locate-and-edit:}} ROME~\cite{rome} locates the factual associations in autoregressive LLMs and conducts accurate and efficient edits by taking MLPs as key-value memories. 
Then, MEMIT~\cite{memit} extends ROME from single-editing to mass-editing.
COMEBA-HK \cite{comeba} identifies the Local Editing Scope and extends MEMIT for sequential editing.
In addition, T-Patcher~\cite{t-patcher} targets the last feed-forward layer of LLMs, adding an additional neuron for each edit. 
\textbf{\textit{Meta learning:}} Recent meta-learning methods use hypernetworks for aiding editing. MEND~\cite{mend} learns a hypernetwork that can decouple the finetuning gradients into the gradient updates that generalize the edits and won't damage the performances on unrelated inputs. To remedy the cancellation effect of MEND, MALMEN~\cite{malmen} uses hypernetwork to produce the weight shifts of editing and formulates the weight shift aggregation as the least square problem. 
\textbf{\textit{Retrieval-based methods:}} Instead of directly editing the model parameters, retrieval-based methods aim to improve the working memory of LLMs to enable model editing. IKE \cite{ike} uses context-edit facts to guide the model when generating edited facts. 
DeCK \cite{deck} employs contrasting knowledge decoding, which enhances the confidence of in-context-based editors in the edited facts. SERAC~\cite{serac} (a modified version dubbed as DEFER~\cite{grace}) records edit items in a file and trains additional scope classifier and counterfactual model to detect, retrieve, and generate the edit-related results. Though the editing retriever and generator are neural networks, they are too small to have the power of LLMs. GRACE~\cite{grace} adopts a discrete codebook of edits for retrieving and replacing the edits' layer representations during inference. From single editing~\cite{rome} to mass editing~\cite{malmen,memit}, and from static editing to sequential~\cite{t-patcher} (continual) or lifelong editing~\cite{grace}, model editing is developing to meet more realistic demands.

\paragraph{Continual Learning}
Continual learning \cite{continual_survey1, continual_survey2} tackles the catastrophic forgetting problem in deep learning models with new knowledge~\cite{de2021continual}, and recent research has focused on various methods in this area. One such method is continual finetuning, where LLMs are refined over time with the arrival of new instances. For instance, a comprehensive study by \cite{lin2022continual} explores continual finetuning extensively. However, it has been observed that regularizing finetuning with continual learning techniques such as Elastic Weight Consolidation~\cite{ewc}, Experience Replay \cite{rolnick2019experience}, and Maximally Interfered Replay \cite{aljundi2019online} can lead to a rapid decay in performance on previous tasks, although it aids in retaining some memory of past inputs. 
This suggests that editing, as opposed to vanilla continual finetuning, presents unique challenges, especially considering that edits are unlikely to be evenly distributed \cite{henn2021principled}. 
One promising direction within the realm of continual learning is the adoption of key-value methods, inspired by advancements in computer vision \cite{liu2021post,van2017neural}. Recent studies have showcased the effectiveness of continual prompt-learning for NLP \cite{wang2022dualprompt,wang2022learning}, particularly in applications like text retrieval \cite{xiong2020approximate}. 
Notably, discrete key-value methods have been shown to excel in handling shifting distributions \cite{trauble2023discrete}, with some recent efforts extending their application to question answering \cite{dai2022lifelong}. 
These methods cache values to ensure that inputs remain within the distribution for downstream encoders, thus facilitating the incorporation of longer-term memory, provided there are adequate computational resources.